\chapter{User Manual}
\label{sec:user_manual}

\section{Introduction}
This user manual provides a practical guide for operating AutoReturn as a unified Gmail- and Slack-based communication workspace. It explains how the application is started, configured, and used in day-to-day workflows.

The focus of this chapter is on the user-facing flow of the system rather than on internal architecture. For that reason, the following sections describe the sequence of use from sign-in and account connection to message handling, settings control, and routine troubleshooting.

\section{Getting Started}
To begin using AutoReturn, the user launches the desktop application, signs in, and connects the supported Gmail and Slack accounts. Once the accounts are active, the main dashboard becomes available for message review, summaries, reply assistance, event inspection, and automation settings.

\begin{enumerate}
    \item Launch AutoReturn and sign in through the authentication interface.
    \item Connect the available Gmail and Slack accounts.
    \item Open the unified inbox dashboard to review synchronized messages.
    \item Select a message to inspect details, summaries, and reply suggestions.
    \item Configure automation policies, quiet rules, or reply preferences from the settings area.
    \item Review event extraction, attachment suggestions, or calendar export options when relevant.
\end{enumerate}

\section{User Interface Guide}
The login and account-connection screens are used to authenticate the user and activate the supported platforms. After authentication, the unified inbox dashboard becomes the central workspace for reviewing Gmail and Slack content, summaries, priority cues, and notification items. Selecting a message opens the detail view, where the user can inspect the content, generate a summary, review a draft reply, or choose a controlled response mode.

The settings area is used to configure quiet hours, allowlists, DND rules, and reply behavior. Event-related workflows are also available when messages contain scheduling details, allowing the user to inspect extracted information and export an ICS file when relevant. Together, these screens represent the main operational path through the application.

\section{Features and Functionality}
AutoReturn combines several user-facing capabilities within one desktop workflow. The main system functions are summarized below:

\begin{enumerate}
    \item \textbf{Unified Communication Workspace:} Gmail and Slack messages are brought into a shared inbox for centralized review.
    \item \textbf{Message Summarization:} Supported messages can be summarized to reduce reading effort and highlight key points quickly.
    \item \textbf{Draft Generation and Reply Assistance:} The system can prepare contextual reply drafts and tone-adjusted suggestions for user review.
    \item \textbf{Priority Classification:} Messages are scored using keyword, deadline, and sender signals to estimate urgency.
    \item \textbf{Automation Policy Control:} Draft-only, plain-reply, auto-reply, allowlist, and DND behavior can be configured by the user.
    \item \textbf{Event Extraction and Calendar Export:} Time-related information can be identified and prepared for calendar-oriented workflows.
    \item \textbf{Voice-Supported Interaction:} The application provides partial support for push-to-talk capture, transcription, and supported command routing.
    \item \textbf{Attachment Assistance:} File requests can be matched against approved locations, with ambiguous cases blocked for manual confirmation.
\end{enumerate}

\section{Frequently Asked Questions (FAQs)}
\begin{enumerate}
    \item \textbf{How do I sign in and connect my accounts?} \\
    Open AutoReturn, complete the authentication step, and connect the available Gmail and Slack accounts from the connection screen. Once the accounts are linked, the system can retrieve communication data and enable the unified workspace.

    \item \textbf{Where can I see all incoming communication?} \\
    All synchronized Gmail and Slack content is shown in the unified inbox dashboard. This is the main area for reviewing recent items, alerts, and prioritized communication.

    \item \textbf{How do I generate a summary or reply draft?} \\
    Select a message from the dashboard and open its detail view. From there, the available assistance features can generate a summary, prepare a draft reply, and present tone-aware suggestions for review.

    \item \textbf{How do I control automation behavior?} \\
    Automation settings are managed through the policy and settings area. Here, the user can configure DND, allowlists, draft-only behavior, plain replies, and other rule-based controls that determine how the system handles communication.
\end{enumerate}
